% Owner: boundary_d. B.4.21 independently refereed by root and limits_ab.
\begin{proposition}[B.4.21 (Proved): a degree-one Kato model with a linear twelve-gon collar]
Start with the $n=1$ Kato family of B.4.10 over a disc
$\Delta_q$, with holomorphic units $a(q),\kappa(q)$ and
$0<|a(q)|<1$.  Pull it back by $q=s^{12}$.
There is a proper modification of this pulled-back threefold,
obtained by six point blow-ups and seventy-two smooth curve
blow-ups supported on its central rational cycle, with the
following properties.  Its total space $\mathcal K^*$ is smooth,
it is proper and flat over $\Delta_s$, and its restriction to
$s\ne0$ is the original pulled-back Kato family.

The strict transform $\mathcal D^*$ of the other Cartier axis
is the smooth semistable Tate surface with multiplier
$\kappa(s^{12})s^{12}$ and central fibre $I_{12}$.  A neighbourhood
of $\mathcal D^*$ in $\mathcal K^*$ is biholomorphic over
$\Delta_s$ to a neighbourhood of the zero section in the
multidegree-zero Tate line with closing factor $a(s^{12})$.
The comparison retains the explicit transverse coordinate
\[
 H=\frac{h}{q},\qquad h=\pi_1^*w-q,
 \tag{B.4.21.1}
\]
on the original punctured family.  Thus its normal comparison,
relative to the original punctured normal frame, has factor
$q^{-1}=s^{-12}$; its inverse has factor $q$.

The persistent elliptic axis has an unchanged neighbourhood.
In particular its normal in the strict transform of the original
central Kato surface still has degree $-1$, rather than $-12$.
The construction supplies an actual linear collar with this
marking.  It does not extend the whole second period, construct
an origin cap for B.4.19, or identify the modified Kato family
with the prescribed torus family. Such a proper filling of
the actual family is excluded by B.4.33; the present
statement retains its auxiliary embedded-collar scope.
\end{proposition}
\begin{proof}
We distinguish the original parameter $q$ from its twelfth root
$s$.  The pulled-back Kato total space is smooth: the original
map to $q$ is a submersion, so locally its pullback has smooth
coordinates consisting of two fibre coordinates and $s$.
The other axis of B.4.20 is Cartier, with closing factor $a(q)$
for its defining function $h$.  At its unique central node its
local equation after this base change is
\[
 h=uv-s^{12}.
 \tag{B.4.21.2}
\]
Away from that point the axis surface is smooth.  The persistent
elliptic axis is disjoint from it throughout the family.

\emph{The six point blow-ups.}
Blow up the singular point of the strict transform of the
Cartier surface, six times in succession.  All these are points
of the smooth ambient threefold.  If its current equation is
$uv=s^m$, with $m\ge2$, the $s$ chart of the ambient point
blow-up has
\[
 u=su_1,\qquad v=sv_1,
 \qquad u_1v_1=s^{m-2}.
 \tag{B.4.21.3}
\]
The $u$ chart has $v=uv_1$, $s=us_1$, and strict-transform
equation
\[
 v_1=u^{m-2}s_1^m;
 \tag{B.4.21.4}
\]
the $v$ chart is symmetric.  The latter charts are smooth.
For $m>2$ the only remaining singular point is the origin in
(B.4.21.3); for $m=2$ that chart has $u_1v_1=1$.
The exponents are $12,10,8,6,4,2$, so precisely six point
blow-ups resolve the surface.

Write $D$ for its smooth strict transform.  Its central fibre is
\[
 F=C_0+C_1+\cdots+C_{11}=\operatorname{div}_D(s).
 \tag{B.4.21.5}
\]
Here $C_0$ is the strict transform of the old nodal rational
curve.  Each of the first five point blow-ups contributes two
rational exceptional curves, and the sixth contributes the
smooth conic in its exceptional plane.  In cyclic order the
two chains are $C_1,\ldots,C_5$ and $C_{11},\ldots,C_7$, with
the conic $C_6$ between them.  Equations (B.4.21.3)--
(B.4.21.4) show that every component has multiplicity one,
that consecutive components meet transversely once, and that
there are no other intersections.  Thus the fibre is $I_{12}$.
Since $F$ is principal, $C_i\cdot F=0$ and $C_i^2=-2$.
The same charts identify $D$ with the minimal regular Tate
model of the given multiplier, including the original closing
unit $\kappa(q)$.

Let $M$ be the pullback to $D$ of the original Cartier normal
line.  It is the Tate line with factor $a(q)$ and multidegree
zero.  A component contracted to the original node has degree
zero for this pullback, and its restriction to $C_0$ is the
pullback of the degree-zero line on the nodal curve.
In the iterated $s$ chart after $j$ point blow-ups the original
function is
\[
 h=s^{2j}(u_jv_j-s^{12-2j}).
\]
The $j$th exceptional ambient divisor therefore has coefficient
$2j$ in the total transform of $\operatorname{div}(h)$.
For $j<6$ its intersection with $D$ consists of
$C_j+C_{12-j}$, and for $j=6$ it is $C_6$.
Consequently the new normal line is
\[
 N_{D/\mathcal K^{(6)}}=M\otimes\mathcal O_D(-Z),
 \quad Z=\sum_{i=0}^{11}z_iC_i,
 \quad z_i=2\min(i,12-i).
 \tag{B.4.21.6}
\]
The list of coefficients is
$(0,2,4,6,8,10,12,10,8,6,4,2)$.
In particular its degree vector is $-4$ on $C_0$, $+4$ on
$C_6$, and zero on the other ten curves.

\emph{The seventy-two curve blow-ups and the actual function.}
For each $i$, blow up the ambient threefold along the current
copy of $C_i$ exactly
\[
 b_i=12-z_i
\]
times, in increasing order of $i$; the entry $b_6=0$ requires
no blow-up.  At every step $C_i$ is a smooth embedded compact
curve.  Its ideal on $D$ is invertible, so the strict transform
of $D$ is canonically isomorphic to $D$ itself.  The normal
line is tensored by $\mathcal O_D(-C_i)$ at that step.
Thus the final normal is
\[
 M\otimes\mathcal O_D(-Z-\textstyle\sum_i b_iC_i)
   =M\otimes\mathcal O_D(-12F)\simeq M.
 \tag{B.4.21.7}
\]
The sum of the $b_i$ is $72$.
The last isomorphism divides the local generator $s^{12}$
of $\mathcal O_D(-12F)$ by $s^{12}$.  Relative to the
original punctured normal identification it is multiplication
by $s^{-12}$, as asserted.

We need more than the line isomorphism (B.4.21.7).
The point-blow-up charts give the following exact local
description near $D$.  At a node $C_i\cap C_{i+1}$ one can
choose ambient coordinates $(x,y,e)$ such that
\[
 D=(e=0),\qquad s=xy,\qquad
 h=e x^{z_i}y^{z_{i+1}}\epsilon,
 \tag{B.4.21.8}
\]
where $\epsilon$ is a holomorphic unit.  At a general point
of $C_i$, omit $y$ and use $s=x$.  These formulas follow
directly from (B.4.21.4): the strict transform is a coordinate
minus a monomial, and the two factors of $s$ are the two
vertical divisors.  They also show transversality of $D$ to
their intersection.  Absorbing a unit in the second base
coordinate preserves this description if different original
coordinate frames are used.

Blowing up $C_i$ is locally the blow-up of $(e,x)$.
In the chart containing the strict transform of $D$, put
$e=xe'$.  Formula (B.4.21.8) keeps $s=xy$ and increases
the exponent of $x$ in $h$ by exactly one.  The old vertical
divisor in this direction no longer meets the strict transform
of $D$.  At an adjacent node the exponent of $y$ is unchanged.
This calculation applies at every repetition and for the
other component as well.  At the end, every vertical ambient
divisor meeting $D$ has $s$-order one and $h$-order twelve.
Hence (B.4.21.8) becomes
\[
 h=s^{12}e\epsilon,
 \qquad H=h/q=e\epsilon.
 \tag{B.4.21.9}
\]
In particular $H$ is holomorphic near all of $D$, has divisor
exactly $D$, and has nonzero normal derivative there.  On the
closing collar its transition is exactly $H\mapsto a(q)H$.
The construction has regularized the actual meromorphic
function, rather than just its normal restriction.

\emph{The relative product charts.}
There is an equivalent useful local target for the preceding
modifications.  The ideal $(h,q)$ is globally defined, since
its changes of generators multiply $h$ by the unit $a(q)$.
Near $D$ the modification maps to its $q$ blow-up chart,
because (B.4.21.9) makes the pullback ideal principal.
At the old node this chart has
\[
 uv=s^{12}(1+H).
\]
Restrict to $|H|<\eta<1$ in a local closing frame and put
$V=v/(1+H)$.  The equation becomes
\[
 uV=s^{12},
\]
times the independent coordinate $H$.
The preceding smooth modification near $D$ is its minimal
resolution times that coordinate.  We verify this assertion
in characters, rather than infer it from the normal line.

Label the exceptional chain so that the orders of $u,V,s$
on $C_i$ are $i,12-i,1$, respectively, with the two end
branches at indices $0,12$.  These are the values in the
explicit surface resolution above.  Along any ambient vertical
divisor meeting $C_i$, the orders of $u$ and $V$ sum to twelve,
because $uV=s^{12}$ and $s$ has order one.  Restriction to
$D$ can only increase those orders.  Since the restricted
orders already sum to twelve, the ambient orders are exactly
$i,12-i$.  No other zero or pole of either character occurs
near $D$, because their product is $s^{12}$.

On the chart of the two adjacent rays $i,i+1$, the functions
\[
 X_i=u/s^i,\qquad Y_i=s^{i+1}/u,
 \qquad X_iY_i=s
 \tag{B.4.21.10}
\]
therefore extend holomorphically.  A meromorphic function on
the smooth ambient space with no codimension-one pole is
holomorphic.  Their restrictions are the two standard minimal
resolution coordinates on $D$.  Together with $H$, whose
normal derivative is nonzero, they have invertible differential
along $D$.  Thus (B.4.21.10) and $H$ are actual biholomorphic
product charts near $D$.  Their character identities give the
standard transitions of the minimal resolution of $A_{11}$.
In particular there is no additional exceptional divisor
meeting $D$ hidden in this product comparison.

\emph{The closing and the uniform linear collar.}
The remaining issue is the global closing, which the Kato
formula determines exactly.  In its blowdown coordinate $x$
and $h=w-q$, that formula is
\[
 (x,h)\longmapsto
 \bigl(\kappa(q)x(a(q)h+q),a(q)h\bigr).
\]
On the regularized collar this reads
\[
 (x,H)\longmapsto
 \bigl(\kappa(q)q x(1+a(q)H),a(q)H\bigr).
 \tag{B.4.21.11}
\]
Shrink the parameter disc so $|a(q)|\le A<1$.
For $|H|<\eta<1$ the product
\[
 G(q,H)=\prod_{j\ge1}(1+a(q)^jH)
 \tag{B.4.21.12}
\]
converges locally uniformly in both variables, is holomorphic
and nonzero, and satisfies
\[
 G(q,a(q)H)=\frac{G(q,H)}{1+a(q)H}.
\]
Uniform convergence follows from
$\sum_{j\ge1}|a(q)^jH|\le\eta A/(1-A)$; the factors have
no zero in the chosen neighbourhood, and the convergent
logarithm series gives a nonzero product and holomorphic inverse.
In the new character $Z=xG(q,H)$, (B.4.21.11) becomes
\[
 (Z,H)\longmapsto
       (\kappa(q)qZ,a(q)H).
 \tag{B.4.21.13}
\]
Multiplication of a toric character by this unit extends on
the product charts (B.4.21.10): it multiplies $X_i$ by the
unit and $Y_i$ by its inverse.  It therefore gives actual
isomorphisms at all twelve nodes as well as on the annular
core.  The exact closing (B.4.21.13) identifies these charts
with the Tate line of factor $a(q)$, over the Tate family
of multiplier $\kappa(q)q$.  Their inverse maps obey the
same identities.  Compactness of the central divisor and
a finite cover by core and node charts allow one common
shrinking of the base and the collar neighbourhood.
This proves the claimed neighbourhood isomorphism.

Every ambient modification used above is a blow-up of a smooth
closed centre and is proper.  The total space remains smooth,
and its map to the parameter disc is flat because a nonzero
base function is a non-zero-divisor in its smooth local rings.
All centres lie over $s=0$, so the punctured family is unchanged.
All centres are disjoint from the persistent elliptic axis,
so its entire neighbourhood is unchanged; in particular its
central normal degree remains $-1$.
We do not assert that the modified ambient map is a submersion
at $s=0$.

The exact exponents, intersection vector, blow-up count and
local substitutions are checked by
\path{code/d_push_degree_one_repair.py} and its saved
output.  The construction stops at the collar.  It does not
extend the second-period action to the complementary Kato
region or remove the exceptional-curve obstruction of B.4.20
to the naive fixed-base plumbing of the entire family.
\end{proof}
